\chapter{The localization}
\label{chap:localization}

Where does the asymmetry between the two directions of a search live? The localization answers
in three parts: \textbf{not} in the dynamics (part 1), \textbf{not} in the state while it is kept
whole (part 2), but \textbf{only} in the erasure of the boundary down to one bit (part 3), where it
is exactly $\cP$ versus $\cNP$. The whole chapter is this project's own; the ingredients that come
from the literature are named in each entry.

\section{Part 1: the dynamics does not carry the arrow}

\subsection{Reversible dynamics}

\begin{definition}[Reversible dynamics]
  \label{def:reversible-dynamics}
  \lean{Sagitta.Arrow.ReversibleDynamics}
  \leanok
  A \emph{reversible dynamics} on a type $\sigma$ is a step $s : \sigma \to \sigma$ with a left
  inverse $b$ (so $b(s(x)) = x$), both of the same cost per step.
  \own{}
\end{definition}

\begin{theorem}[Going back costs the same as going forward]
  \label{thm:back-iterate}
  \lean{Sagitta.Arrow.ReversibleDynamics.back_iterate_step_iterate, Sagitta.Arrow.ReversibleDynamics.costBack_eq_costForward, Sagitta.Arrow.ReversibleDynamics.step_injective}
  \leanok
  \uses{def:reversible-dynamics}
  For every $n$ and $x$, $b^n(s^n(x)) = x$, and the cost of $n$ steps back equals the cost of $n$
  steps forward. The step is injective.
  \own{}
\end{theorem}

\begin{proof}
  \leanok
  Induction on $n$.
\end{proof}

\begin{theorem}[Any dynamics becomes reversible with a history]
  \label{thm:history-dynamics}
  \lean{Sagitta.Arrow.historyDynamics, Sagitta.Arrow.history_simulates, Sagitta.Arrow.history_recovers, Sagitta.Arrow.injective_iff_exists_back}
  \leanok
  \uses{def:reversible-dynamics, thm:back-iterate}
  For any $f : \sigma \to \sigma$, the dynamics on $\sigma \times \sigma^*$ that applies $f$ and pushes
  the old state on a history list is reversible, simulates $f$, and recovers the initial state after
  $n$ steps back. A map has a left inverse if and only if it is injective.
  \source{this is the abstract shadow of Bennett's construction \cite{Bennett73}; the
  Turing-machine version is Chapter~\ref{chap:bennett}.}
\end{theorem}

\begin{proof}
  \leanok
  Pop the history.
\end{proof}

\subsection{Time reversal: information is conserved (discrete Noether)}

\begin{definition}[Koopman operator]
  \label{def:koopman}
  \lean{Sagitta.Arrow.Dynamics.koopman, Sagitta.Arrow.Dynamics.mult, Sagitta.Arrow.Dynamics.phase, Sagitta.Arrow.Dynamics.norm2}
  \leanok
  For $f : \sigma \to \sigma$ on a finite set, the Koopman operator is the linear map
  $K_f \varphi = \varphi \circ f$ on $\mathbb{C}^\sigma$; $M_\varphi$ is multiplication by $\varphi$,
  $e^{i\theta Q}$ the phase of a real quantity $Q$, and $\|\varphi\|^2 = \sum_x |\varphi(x)|^2$.
  \source{\cite{Koopman31}.}
\end{definition}

\begin{theorem}[Discrete Noether]
  \label{thm:noether}
  \lean{Sagitta.Arrow.Dynamics.conserved_iff_const_on_orbits, Sagitta.Arrow.Dynamics.koopman_noether}
  \leanok
  \uses{def:koopman}
  For a permutation $f$, a quantity $Q$ is conserved ($Q \circ f = Q$) if and only if it is constant
  on the orbits of the group generated by $f$. For any $f$ on a finite set, $Q \circ f = Q$ if and
  only if $K_f$ commutes with $M_{e^{i\theta Q}}$ for every $\theta$.
  \own{} (a discrete analogue of Noether's theorem in Koopman's picture).
\end{theorem}

\begin{proof}
  \leanok
  Orbits: induction on the powers of $f$. Koopman: evaluate the commutator on indicator functions and
  separate the values of $Q$ with a suitable $\theta$.
\end{proof}

\begin{theorem}[Unitarity is bijectivity]
  \label{thm:koopman-isometry}
  \lean{Sagitta.Arrow.Dynamics.koopman_isometry_iff, Sagitta.Arrow.Dynamics.erasure_not_isometry}
  \leanok
  \uses{def:koopman}
  $K_f$ preserves $\|\cdot\|^2$ if and only if $f$ is bijective. Erasure (the constant map on
  $\{0,1\}$) is not an isometry.
  \own{}
\end{theorem}

\begin{proof}
  \leanok
  Evaluate the norm on indicator functions of single states and of sets.
\end{proof}

\begin{theorem}[Reversible dynamics neither creates nor destroys information]
  \label{thm:card-image}
  \lean{Sagitta.Arrow.ReversibleDynamics.card_image_iterate, Sagitta.Arrow.ReversibleDynamics.koopman_isometry, Sagitta.Arrow.Dynamics.entropy_pushforward}
  \leanok
  \uses{def:reversible-dynamics, thm:koopman-isometry}
  A reversible dynamics keeps the size of every finite set of states after $n$ steps, its Koopman
  operator is an isometry on a finite type, and a permutation preserves the Shannon entropy of a
  distribution.
  \own{}
\end{theorem}

\begin{proof}
  \leanok
  The step is injective (Theorem~\ref{thm:back-iterate}).
\end{proof}

\subsection{Time translation: the system does not know the time}

\begin{definition}[Homogeneous automaton]
  \label{def:automaton}
  \lean{Sagitta.Arrow.Dynamics.Automaton, Sagitta.Arrow.Dynamics.Automaton.run, Sagitta.Arrow.Dynamics.Automaton.step}
  \leanok
  An automaton that consults a verifier has states $\sigma$, an initial state, a query $q(x)$ and an
  output $o(x) \in Y \cup \{\bot\}$ for each state, and a transition $\delta(x, r)$ on the answer
  $r$. It is \emph{homogeneous in time}: there is no step counter, and what it knows about the time
  is what it keeps in $\sigma$.
  \own{}
\end{definition}

\begin{theorem}[The Noether charge of time translation]
  \label{thm:noether-charge}
  \lean{Sagitta.Arrow.Dynamics.noether_charge, Sagitta.Arrow.Dynamics.does_not_know_the_time}
  \leanok
  \uses{def:automaton}
  Two runs against the same verifier that reach the same state, one at time $i$ and the other at
  time $j$, without having stopped, continue identically. In particular a run that comes back to a
  state continues from $i$ and from $j$ in the same way: it does not know the time.
  \own{}
\end{theorem}

\begin{proof}
  \leanok
  Split the run at $i$ and at $j$.
\end{proof}

\begin{definition}[Clocked automaton]
  \label{def:clocked}
  \lean{Sagitta.Arrow.Dynamics.ClockedAutomaton, Sagitta.Arrow.Dynamics.ClockedAutomaton.run, Sagitta.Arrow.Dynamics.ClockedAutomaton.homogenize}
  \leanok
  \uses{def:automaton}
  A clocked automaton also reads a time register $\kappa$ that advances on its own. Its
  \emph{homogenization} is the homogeneous automaton on $\sigma \times \kappa$.
  \own{}
\end{definition}

\begin{theorem}[A clock is paid in memory]
  \label{thm:homogenize}
  \lean{Sagitta.Arrow.Dynamics.ClockedAutomaton.homogenize_run, Sagitta.Arrow.Dynamics.card_clock_bits}
  \leanok
  \uses{def:clocked}
  The homogenization has the same runs as the clocked automaton, and a clock of $c$ bits multiplies
  the number of states by $2^c$.
  \own{}
\end{theorem}

\begin{proof}
  \leanok
  Induction on the run.
\end{proof}

\begin{theorem}[The cube law]
  \label{thm:cube-law}
  \lean{Sagitta.Arrow.Dynamics.cube_law, Sagitta.Arrow.Dynamics.cube_law_general, Sagitta.Arrow.Dynamics.cube_law_clock}
  \leanok
  \uses{def:automaton, thm:homogenize}
  If automata indexed by at most $2^L$ programs and $2^s$ advice strings, with at most $2^S$ states
  each, find any of $N$ targets with $t$ yes/no queries, then
  $N \le 2^L \cdot 2^s \cdot \min(t+1, 2^{S-1})$. With answers of $b$ bits the bound is
  $2^L \cdot 2^s \cdot \min(2^{bt}, 2^S)$; with a clock of $c$ bits, $S$ becomes $S + c$.
  \own{}
\end{theorem}

\begin{proof}
  \leanok
  Count the targets found along the all-``no'' trajectory, bounded both by the time and by the
  number of states, and take a union over programs and advice.
\end{proof}

\begin{theorem}[Halt or cycle; the CHT with its costs]
  \label{thm:cht}
  \lean{Sagitta.Arrow.Dynamics.halt_or_cycle, Sagitta.Arrow.Dynamics.halts_before_two_pow, Sagitta.Arrow.Dynamics.cht_time, Sagitta.Arrow.Dynamics.clock_in_memory}
  \leanok
  \uses{thm:noether-charge, thm:homogenize}
  A deterministic system on a finite set either halts within $|\sigma|$ steps or never halts and is
  eventually periodic. Hence a homogeneous automaton with $|\sigma| \le 2^S$ that answers, answers
  in fewer than $2^S$ queries; with a clock of $c$ bits, in fewer than $2^{S+c}$.
  \own{} (``CHT'' is this project's name for the bound; the pigeonhole argument is classical).
\end{theorem}

\begin{proof}
  \leanok
  Pigeonhole on the trace, and the Noether charge after a repeated state.
\end{proof}

\begin{theorem}[The CHT bound is exponential, and attained]
  \label{thm:cht-exponential}
  \lean{Sagitta.Arrow.Dynamics.counter, Sagitta.Arrow.Dynamics.cht_exponential_not_polynomial}
  \leanok
  \uses{thm:cht}
  For every $n \ge 1$ there is a homogeneous automaton with $n+1$ bits of memory (a counter) that
  always answers within the budget $2^{n+1}$ but does not answer in $2^n - 1$ queries.
  \own{}
\end{theorem}

\begin{proof}
  \leanok
  The counter walks through $2^n$ states before it may answer.
\end{proof}

\subsection{Part 1 for Turing machines}

\begin{definition}[A machine with bounded space is a homogeneous automaton]
  \label{def:toAutomaton}
  \lean{Sagitta.Arrow.TM.toAutomaton, Sagitta.Arrow.TM.spaceBits}
  \leanok
  \uses{def:automaton, def:cslib-tm, def:cslib-time-space, def:cslib-configbound}
  A cslib multi-tape machine with $k$ work tapes, run on an input of length $n$ in space $S$, is a
  homogeneous automaton whose states are its configurations without the output (the state, the
  input head and the windows of the work tapes), at most $2^{\mathrm{spaceBits}}$ of them.
  \own{}
\end{definition}

\begin{theorem}[cslib's machines do not know the time]
  \label{thm:tm-dynamics}
  \lean{Sagitta.Arrow.TM.run_toAutomaton, Sagitta.Arrow.TM.card_states_bits, Sagitta.Arrow.TM.does_not_know_the_time_tm, Sagitta.Arrow.TM.cube_law_tm, Sagitta.Arrow.TM.halt_or_cycle_tm, Sagitta.Arrow.TM.halts_soon_tm}
  \leanok
  \uses{def:toAutomaton, thm:noether-charge, thm:cube-law, thm:cht}
  A machine that runs in space $S$ on an input does not know the time (if its configuration repeats
  before it halts, the run repeats), obeys the cube law, halts or cycles, and if it halts, halts
  within $2^{\mathrm{spaceBits}}$ steps.
  \own{}
\end{theorem}

\begin{proof}
  \leanok
  The automaton of Definition~\ref{def:toAutomaton} runs in lockstep with the machine.
\end{proof}

\begin{theorem}[CHT tower, piece 1: the $\cPSPACE$ bound]
  \label{thm:cht-pspace}
  \lean{Sagitta.Arrow.TM.cht_pspace}
  \leanok
  \uses{thm:tm-dynamics}
  A machine that runs in polynomial space and halts on an input halts within $2^{b(n)}$ steps for a
  polynomially bounded $b$.
  \source{the classical bound $\mathsf{SPACE}(s) \subseteq \TIME(2^{O(s)})$; \cite{AroraBarak},
  Chapter~4.}
\end{theorem}

\begin{proof}
  \leanok
  Theorem~\ref{thm:tm-dynamics}, with $\mathrm{spaceBits}$ polynomial in $n$.
\end{proof}

\begin{theorem}[Localization, part 1, complete]
  \label{thm:localization-one}
  \lean{Sagitta.Arrow.Dynamics.localization_one, Sagitta.Arrow.TM.localization_one_tm}
  \leanok
  \uses{thm:noether, thm:koopman-isometry, thm:card-image, thm:noether-charge, thm:homogenize, thm:cht, thm:arrow-independent, thm:hd-iff-arrow, thm:tm-dynamics, thm:cht-pspace, thm:cht-exponential, thm:hd-iff-arrow-tm}
  Neither symmetry of time decides the arrow. Time reversal conserves information (discrete Noether,
  unitarity iff bijectivity, sizes of sets kept); time translation makes a homogeneous system blind
  to the clock and charges the clock to memory, and the CHT bounds by $2^S$ the time of a system of
  $S$ bits. These facts hold in every model with no hypothesis, while the arrow holds in one model
  and fails in another; what decides it is the decision hypothesis, which is about the bit, not the
  dynamics. The same holds for cslib's machines.
  \own{}
  \fidelity{a conjunction of the proved facts. ``Neither symmetry decides the arrow'' is the essay's
  reading, not a formal statement; \texttt{arrow\_independent} uses models without dynamics.}
\end{theorem}

\begin{proof}
  \leanok
  The conjunction of the cited results.
\end{proof}

\section{Part 2: the full boundary, and the interpolation}

\begin{definition}[Completing a prefix]
  \label{def:complete}
  \lean{Sagitta.Arrow.complete}
  \leanok
  Given a verifier $V$ on bit strings and a prefix $u$, $\mathrm{complete}(V, u, k)$ searches the
  $2^k$ completions of $u$ by depth-first search and returns the first $w$ with $V(w) = 1$, together
  with the number of calls to $V$.
  \own{}
\end{definition}

\begin{theorem}[The full boundary]
  \label{thm:full-boundary}
  \lean{Sagitta.Arrow.full_boundary}
  \leanok
  \uses{def:complete}
  When the whole witness $w$ is kept ($V(w) = 1$), going back is one verification:
  $\mathrm{complete}(V, w, 0) = (w, 1)$.
  \own{}
\end{theorem}

\begin{proof}
  \leanok
  By definition.
\end{proof}

\begin{theorem}[The interpolation]
  \label{thm:interpolation}
  \lean{Sagitta.Arrow.interpolation, Sagitta.Arrow.complete_calls_le, Sagitta.Arrow.complete_sound, Sagitta.Arrow.complete_isSome}
  \leanok
  \uses{def:complete}
  With $m$ of the $n$ bits of a witness kept (a prefix $u$ of length $m$ that has a good completion
  of length $n$), going back finds a witness of length $n$ extending $u$ with at most $2^{n-m}$ calls
  to the verifier, so at cost at most $2^{n-m} c$ when one call costs $c$.
  \own{}
\end{theorem}

\begin{proof}
  \leanok
  Induction on the number of missing bits: each level doubles the calls at most.
\end{proof}

\begin{theorem}[The interpolation is tight in a black box]
  \label{thm:boundary-lower-bound}
  \lean{Sagitta.Arrow.boundary_lower_bound}
  \leanok
  \uses{lem:adversary}
  Any query tree over the $2^k$ completions that decides whether one of them is good asks, on the
  all-zero oracle, about all $2^k$.
  \own{}
\end{theorem}

\begin{proof}
  \leanok
  \uses{lem:adversary}
  The adversary lemma with $|\alpha| = 2^k$.
\end{proof}

\section{Part 3: one bit}

\begin{theorem}[With one bit, going back is cheap exactly when $\cP = \cNP$]
  \label{thm:searchEasy-iff}
  \lean{Sagitta.Arrow.Model.searchEasy_iff_peqNP}
  \leanok
  \uses{def:model-classes}
  In every model, search is easy if and only if $\cP = \cNP$.
  \source{search-to-decision \cite{Goldreich}, Chapter~2; here derived from the closure properties
  alone.}
\end{theorem}

\begin{proof}
  \leanok
  Decision reduces to search through the domain; search reduces to decision through the prefix
  problem, one call per bit.
\end{proof}

\begin{theorem}[One bit, for Turing machines]
  \label{thm:oneBit}
  \lean{Sagitta.Arrow.TM.oneBit, Sagitta.Arrow.TM.searchEasy_iff}
  \leanok
  \uses{thm:cslib-search-decision, def:tmModel}
  $\cFNP \subseteq \cFP$ for every encoding if and only if $\cNP \subseteq \cP$ for every encoding.
\end{theorem}

\begin{proof}
  \leanok
  \uses{thm:cslib-search-decision}
  It is cslib-frontier's theorem (Theorem~\ref{thm:cslib-search-decision}).
\end{proof}

\section{The localization theorem}

\begin{theorem}[Localization]
  \label{thm:localization}
  \lean{Sagitta.Arrow.localization, Sagitta.Arrow.TM.localizationTM}
  \leanok
  \uses{thm:back-iterate, thm:full-boundary, thm:interpolation, thm:searchEasy-iff, thm:oneBit}
  In every model, and for Turing machines:
  \begin{enumerate}
    \item a reversible dynamics, run back from its final state, recovers the initial state at the
      same cost as running it forward;
    \item with the full boundary going back is one verification, and with $m$ of the $n$ bits it is at
      most $2^{n-m}$ verifications;
    \item with one bit (only the answer ``yes'') going back is cheap for every search problem if and
      only if $\cP = \cNP$.
  \end{enumerate}
  \own{}
  \fidelity{a conjunction of the three parts. The word ``only'' is the essay's reading, not a formal
  statement. In the abstract core the equality of the cost back and forward is definitional; for
  Turing machines there is no general costed inverse.}
\end{theorem}

\begin{proof}
  \leanok
  The conjunction of the three parts.
\end{proof}
